\chapter{High-Frequency Econometrics}\label{ch:qm:high-frequency-econometrics}

Measured from one-second changes in the last traded price, a stock's annualised volatility is 60\%; from five-minute changes it is 25\%. The
difference is not the market's. The trades alternate between the bid and the ask, and the prices sit on a one-cent grid; at one second those two
effects are most of the measured variation, and at five minutes almost none of it. Chapter~18 promised that \omterm{def:qm:volatility-models:rv}{realised variance} converges to the true
\omterm{def:qm:volatility-models:rv}{integrated variance} as the sampling gets finer; this chapter is about why, in real prices, it does the opposite, and what to do about it: the noise
and the \omterm{def:qm:high-frequency-econometrics:signature}{signature plot} that reveals it, the optimal sampling interval, \omterm{def:qm:estimation:estimator}{estimators} that remove the noise instead of avoiding it, the correlation that
vanishes when two assets trade at different moments, and jumps.

The chapter's market is simulated, so that every estimate can be checked against the truth: an efficient log price with 25\% annual volatility, around
36 dollars, traded once a second at the bid or the ask of a one-cent grid (the bid is the efficient price rounded down to the cent, the ask one cent
higher), over 6.5-hour days of 23\,400 seconds.

\section{Microstructure noise and the signature plot}

\begin{definition}[Microstructure noise, bid--ask bounce]\label{def:qm:high-frequency-econometrics:noise}
\emph{Microstructure noise}\index{microstructure noise} is the difference between an observed price and the efficient price, created by the trading process: the
spread, the price grid, the discreteness and staleness of quotes. The \emph{bid--ask bounce}\index{bid--ask bounce} is its best-known form: successive trades at the
bid and at the ask make observed returns alternate in sign when the efficient price has not moved.
\end{definition}

If the observed log price is $Y_i = X_i + u_i$ with iid noise $u_i$ of variance $\omega^2$ independent of the efficient $X$, each observed return carries $u_i - u_{i-1}$, so
\omterm{def:qm:volatility-models:rv}{realised variance} from $n$ returns has
\[
\E[\mathrm{RV}_n] = \mathrm{IV} + 2n\omega^2:
\]
a bias that grows linearly with the sampling frequency and swamps the \omterm{def:qm:volatility-models:rv}{integrated variance} at high frequency. The noise also makes consecutive returns
negatively correlated, with first autocovariance $-\omega^2$.

\begin{definition}[Roll's estimator]\label{def:qm:high-frequency-econometrics:roll}
\emph{Roll's estimator}\index{Roll's estimator} (Roll, 1984) of the effective bid--ask spread is $2\sqrt{-\gamma_1}$, with $\gamma_1$ the first autocovariance of trade-price
changes: under pure bounce around an efficient price, trades are at mid $\pm c/2$ with independent signs, and $\gamma_1 = -c^2/4$.
\end{definition}

On the simulated stock, whose spread is one cent, \omterm{def:qm:high-frequency-econometrics:roll}{Roll's estimator} averages 1.14 cents over thirty days: close, and biased up because the efficient price also
moves the bid within the grid.

\begin{definition}[Signature plot]\label{def:qm:high-frequency-econometrics:signature}
A \emph{signature plot}\index{signature plot} draws average realised volatility against the sampling interval: flat
without noise, rising toward the shortest intervals with it (Andersen, Bollerslev, Diebold and Labys, 2000, who named it the volatility signature plot).
\end{definition}

\Cref{fig:qm:high-frequency-econometrics:signature} is the chapter's \omterm{def:qm:high-frequency-econometrics:signature}{signature plot}. From trades, realised volatility is 60.1\% at one second, 35.1\% at five
seconds, 26.1\% at one minute and 25.4\% at five minutes; from mid quotes, which do not bounce but still sit on the grid, 36.8\% at one second; from the
efficient price, 25.0\% at every interval.

\begin{omfigure}
\begin{tikzpicture}
\begin{semilogxaxis}[width=10cm,height=5cm,xlabel={sampling interval (seconds)},ylabel={annualised volatility (\%)},
  tick label style={font=\small},label style={font=\small},xmin=1,xmax=1800,ymin=20,ymax=65,grid=major,grid style={gray!25},
  xtick={1,10,60,300,1800},xticklabels={1,10,60,300,1800},
  legend style={font=\footnotesize,at={(0.97,0.97)},anchor=north east,fill=white,draw=none},table/col sep=comma]
\addplot[omThm,very thick,mark=*,mark size=1.3pt] table[x=seconds,y=trade]{figdata/methods/21-high-frequency-econometrics/signature.csv};
\addplot[black,thick,mark=square*,mark size=1.3pt] table[x=seconds,y=mid]{figdata/methods/21-high-frequency-econometrics/signature.csv};
\addplot[omDef,very thick,dashed] table[x=seconds,y=efficient]{figdata/methods/21-high-frequency-econometrics/signature.csv};
\legend{trade prices,mid quotes,efficient price}
\end{semilogxaxis}
\end{tikzpicture}

\omcaption{\omterm{def:qm:high-frequency-econometrics:signature}{Signature plot} of the simulated stock (efficient volatility 25\%, one-cent grid around 36 dollars, a trade every second at the bid or the ask), averaged
over thirty days: realised volatility against the sampling interval. The bounce and the grid add 35 points at one second and almost nothing beyond two minutes. Data:
the chapter's tutorial, seeded.}\label{fig:qm:high-frequency-econometrics:signature}
\end{omfigure}

The noise variance can be estimated from the highest frequency itself, $\hat\omega^2 = \mathrm{RV}_n/(2n)$ (which also counts $\mathrm{IV}/2n$, negligible here): 1.75
basis points of standard deviation, against the true 1.60. Sampling less often trades bias for variance, and the trade has an optimum.

\begin{proposition}[Optimal sampling under iid noise]\label{prop:qm:high-frequency-econometrics:optimal}
With iid noise of variance $\omega^2$ and constant volatility, the mean squared error of $\mathrm{RV}_n$ is approximately $2\,\mathrm{IQ}/n + (2n\omega^2)^2$, where $\mathrm{IQ} =
\int\sigma^4$ is the integrated quarticity; it is minimised at $n^* = (\mathrm{IQ}/(4\omega^4))^{1/3}$ (Bandi and Russell, 2008).
\end{proposition}

\begin{proof}
The discretisation error of \omterm{def:qm:volatility-models:rv}{realised variance} has variance $2\,\mathrm{IQ}/n$ (chapter~18's $\sqrt{2/m}$ with constant volatility), and the noise adds the bias $2n\omega^2$
(plus variance terms of lower order at moderate $n$). Setting the derivative $-2\,\mathrm{IQ}/n^2 + 8n\omega^4$ to zero gives $n^{*3} = \mathrm{IQ}/(4\omega^4)$.
\end{proof}

With the estimated noise, the noise-to-signal ratio $\omega^2/\mathrm{IV}$ is $1.24 \times 10^{-4}$ and $n^* = 254$ returns a day, one every 92 seconds (82 with the true
noise). The simulation agrees: the root-mean-squared relative error of daily \omterm{def:qm:volatility-models:rv}{realised variance} is 4.80 at one second, 0.18 at 30 seconds, 0.11 at 60, 0.12 at 120
and 0.14 at five minutes (\cref{fig:qm:high-frequency-econometrics:mse}). One-second sampling overstates the variance 5.8-fold; the optimum avoids that bias at the
cost of using a fraction of the data.

\begin{omfigure}
\begin{tikzpicture}
\begin{loglogaxis}[width=10cm,height=5cm,xlabel={sampling interval (seconds)},ylabel={RMSE of RV / IV},
  tick label style={font=\small},label style={font=\small},xmin=1,xmax=1800,ymin=0.08,ymax=6,grid=major,grid style={gray!25},
  xtick={1,10,60,300,1800},xticklabels={1,10,60,300,1800},ytick={0.1,0.2,0.5,1,2,5},yticklabels={0.1,0.2,0.5,1,2,5},table/col sep=comma]
\addplot[omDef,very thick,mark=*,mark size=1.5pt] table[x=seconds,y=rmse]{figdata/methods/21-high-frequency-econometrics/mse.csv};
\addplot[omThm,thick,dashed,sharp plot] coordinates {(92,0.08) (92,6)};
\end{loglogaxis}
\end{tikzpicture}

\omcaption{Root-mean-squared relative error of daily \omterm{def:qm:volatility-models:rv}{realised variance} from trade prices against the sampling interval, over sixty simulated days. The dashed line
is the Bandi--Russell optimum of 92 seconds computed from the estimated noise variance. Data: the chapter's tutorial, seeded.}\label{fig:qm:high-frequency-econometrics:mse}
\end{omfigure}

\section{Estimators robust to noise}

Sparse sampling throws away most of the data. Three families use it all and remove the noise.

\begin{definition}[Two-scales realised variance]\label{def:qm:high-frequency-econometrics:tsrv}
The \emph{two-scales realised variance}\index{two-scales realised variance} (Zhang, Mykland and Aït-Sahalia, 2005) averages the \omterm{def:qm:volatility-models:rv}{realised variances} of the $K$ offset sparse grids
and subtracts the noise bias estimated from the full grid: $\mathrm{TSRV} = \overline{\mathrm{RV}}^{(K)} - (\bar n/n)\mathrm{RV}^{(\mathrm{all})}$, $\bar n = (n - K + 1)/K$.
\end{definition}

Averaging over the $K$ grids reduces the discretisation error of the sparse \omterm{def:qm:estimation:estimator}{estimator}; the subtraction removes its remaining bias $2\bar n\omega^2$. It is consistent under
iid noise, at the slow rate $n^{-1/6}$.

\begin{definition}[Realised kernel, pre-averaging estimator]\label{def:qm:high-frequency-econometrics:kernel}
A \emph{realised kernel}\index{realised kernel} (Barndorff-Nielsen, Hansen, Lunde and Shephard, 2008) adds weighted autocovariances of the high-frequency returns to their
variance, $\gamma_0 + \sum_{h=1}^Hk\bigl(\frac{h-1}H\bigr)(\gamma_h + \gamma_{-h})$, with a smooth weight such as Parzen's, to cancel the negative autocovariance the noise creates.
The \emph{pre-averaging estimator}\index{pre-averaging estimator} (Jacod, Li, Mykland, Podolskij and Vetter, 2009) averages returns over short overlapping windows with a
weight function, which shrinks the noise in each window, then squares and rescales the averages and subtracts the remaining bias.
\end{definition}

The kernel is the long-run variance of chapter~11 applied within a day: the noise makes high-frequency returns an MA(1), and a HAC-type sum recovers the variance of the
efficient part. Over thirty simulated days, the average annualised volatility from trades is 25.4\% by five-minute RV, 25.2\% by two-scales RV ($K = 300$), 25.8\% by RV at the
optimal interval, 25.1\% by pre-averaging (windows of 150 seconds) and 25.1\% by the \omterm{def:qm:high-frequency-econometrics:kernel}{realised kernel} ($H = 60$). The day-to-day relative standard deviations separate them
(\cref{fig:qm:high-frequency-econometrics:estimators}): 14.3\% for five-minute RV, 11.9\% for two-scales, 7.9\% at the optimal interval, 7.2\% for pre-averaging and 4.8\% for the
kernel.

\begin{omfigure}
\begin{tikzpicture}
\begin{axis}[width=9cm,height=4.8cm,xbar,bar width=9pt,
  xlabel={relative standard deviation across days (\%)},
  ytick={0,1,2,3,4},yticklabels={5-minute RV,two-scales RV,RV at 92 s,pre-averaging,realised kernel},
  tick label style={font=\small},label style={font=\small},
  xmin=0,xmax=16,ymin=-0.6,ymax=4.6,grid=major,grid style={gray!25},table/col sep=comma]
\addplot[fill=omDef!60,draw=omDef,nodes near coords,nodes near coords style={font=\footnotesize,anchor=west,xshift=1pt}]
  table[x=sd,y=k]{figdata/methods/21-high-frequency-econometrics/estimators.csv};
\end{axis}
\end{tikzpicture}

\omcaption{Precision of five daily volatility \omterm{def:qm:estimation:estimator}{estimators} on the simulated stock's trade prices, thirty days: the standard deviation of the daily estimate relative to the true
\omterm{def:qm:volatility-models:rv}{integrated variance}. All are nearly unbiased (average volatilities between 25.1\% and 25.8\% for a true 25\%). Data: the chapter's tutorial,
seeded.}\label{fig:qm:high-frequency-econometrics:estimators}
\end{omfigure}

\section{Asynchronous observation and the Epps effect}

\begin{definition}[Non-synchronous trading, Epps effect]\label{def:qm:high-frequency-econometrics:epps}
\emph{Non-synchronous trading}\index{non-synchronous trading} is the observation of different assets at different, random times. The \emph{Epps
effect}\index{Epps effect} (Epps, 1979) is the resulting fall of realised correlation toward zero as the sampling interval shortens: with previous-tick sampling on a fine grid,
most intervals contain a price change of one asset and none of the other.
\end{definition}

Two simulated stocks with a true correlation of 0.6, trading at Poisson times every 5 and every 15 seconds on average, show it cleanly
(\cref{fig:qm:high-frequency-econometrics:epps}): realised correlation from previous-tick prices is 0.03 at one second, 0.21 at ten seconds, 0.49 at one minute, 0.59 at five
minutes and 0.63 at thirty minutes (where the sampling error of few returns takes over).

\begin{definition}[Hayashi--Yoshida estimator, refresh-time sampling]\label{def:qm:high-frequency-econometrics:hy}
The \emph{Hayashi--Yoshida estimator}\index{Hayashi--Yoshida estimator} (Hayashi and Yoshida, 2005) of the covariation sums the products of all pairs of returns, one from each
asset, whose observation intervals overlap, with no interpolation. \emph{Refresh-time sampling}\index{refresh-time sampling} samples both assets at the first time each has traded
since the previous sampling time.
\end{definition}

The \omterm{def:qm:high-frequency-econometrics:hy}{Hayashi--Yoshida estimator} is unbiased for the covariation of the efficient prices without noise, whatever the observation times, because every piece of the covariance
is counted once. On the two stocks it gives an average correlation of 0.594 (true 0.6), with a day-to-day standard deviation of 0.030. \omterm{def:qm:high-frequency-econometrics:hy}{Refresh-time sampling} alone gives 0.48:
it synchronises the sampling, but the last prices it pairs were still observed at different moments.

\begin{omfigure}
\begin{tikzpicture}
\begin{semilogxaxis}[width=10cm,height=5cm,xlabel={sampling interval (seconds)},ylabel={realised correlation},
  tick label style={font=\small},label style={font=\small},xmin=1,xmax=1800,ymin=0,ymax=0.75,grid=major,grid style={gray!25},
  xtick={1,10,60,300,1800},xticklabels={1,10,60,300,1800},
  legend columns=3,legend style={font=\footnotesize,at={(0.5,-0.3)},anchor=north,draw=none,fill=none,
  /tikz/every even column/.append style={column sep=6pt}},table/col sep=comma]
\addplot[omDef,very thick,mark=*,mark size=1.5pt] table[x=seconds,y=corr]{figdata/methods/21-high-frequency-econometrics/epps.csv};
\addplot[omThm,very thick,dashed] table[x=seconds,y=hy]{figdata/methods/21-high-frequency-econometrics/epps.csv};
\addplot[black,thin] table[x=seconds,y=truth]{figdata/methods/21-high-frequency-econometrics/epps.csv};
\legend{previous-tick,Hayashi--Yoshida,true correlation}
\end{semilogxaxis}
\end{tikzpicture}

\omcaption{The \omterm{def:qm:high-frequency-econometrics:epps}{Epps effect}: realised correlation of two simulated stocks (true correlation 0.6) that trade at independent Poisson times, on average every 5 and 15 seconds,
against the sampling interval, averaged over forty days; and the Hayashi--Yoshida estimate from all the trades (0.594). Data: the chapter's tutorial,
seeded.}\label{fig:qm:high-frequency-econometrics:epps}
\end{omfigure}

\section{Jumps at high frequency}

\begin{definition}[Bipower variation]\label{def:qm:high-frequency-econometrics:bipower}
\emph{Bipower variation}\index{bipower variation} (Barndorff-Nielsen and Shephard, 2004) is $\frac\pi2\sum_i|r_i||r_{i-1}|$: products of adjacent absolute returns, scaled by
$1/\E|Z|^2 = \pi/2$.
\end{definition}

\omterm{def:qm:volatility-models:rv}{Realised variance} converges to the whole \omterm{def:qm:brownian-motion:qv}{quadratic variation}, the \omterm{def:qm:volatility-models:rv}{integrated variance} plus the sum of squared jumps; \omterm{def:qm:high-frequency-econometrics:bipower}{bipower variation} converges to the \omterm{def:qm:volatility-models:rv}{integrated variance}
alone, because a jump enters only through products with a neighbouring, diffusive return, which vanish as the sampling gets finer. The difference estimates the jump part. At
five-minute sampling, over forty simulated days with a 1\% jump at midday (29\% of that day's \omterm{def:qm:brownian-motion:qv}{quadratic variation}), \omterm{def:qm:volatility-models:rv}{realised variance} averages 1.43 times the \omterm{def:qm:volatility-models:rv}{integrated
variance} and \omterm{def:qm:high-frequency-econometrics:bipower}{bipower variation} 1.16; without the jump, 1.00 and 1.00. At five minutes the jump still leaks into \omterm{def:qm:high-frequency-econometrics:bipower}{bipower variation} through its two neighbouring products; the
cure is finer sampling, which brings back the noise, and so jump detection and noise-robust estimation have to be combined in practice.

\section{Tutorial: the volatility that depended on the clock}\label{tut:qm:high-frequency-econometrics:clock}

\begin{tutorial}
\textbf{Goal.} Measure a simulated stock's volatility at every sampling frequency, find the optimal interval, and remove the noise instead of avoiding it.
\textbf{End state:} \cref{fig:qm:high-frequency-econometrics:signature,fig:qm:high-frequency-econometrics:mse,fig:qm:high-frequency-econometrics:epps}, the optimal interval of 92
seconds, and the kernel's 4.8\% daily error.
\begin{tutsteps}
\item \textbf{\omterm{def:qm:high-frequency-econometrics:tsrv}{Two-scales realised variance} and the \omterm{def:qm:high-frequency-econometrics:kernel}{realised kernel}.}
\omcode{code/firm/hfvol/firm_hfvol.py}{44}{64}{Two-scales realised variance and the Parzen realised kernel.}\label{lst:qm:high-frequency-econometrics:tsrv}
\item \textbf{The Hayashi--Yoshida covariance} without interpolation.
\omcode{code/firm/hfvol/firm_hfvol.py}{100}{110}{The Hayashi--Yoshida estimator.}\label{lst:qm:high-frequency-econometrics:hy}
\item \textbf{Run} \texttt{signature()}, \texttt{optimal\_interval()}, \texttt{estimators()}, \texttt{epps()} and \texttt{jumps()} in \texttt{qm\_hf.py}, then
\texttt{fig\_hf.py}.
\end{tutsteps}
\textbf{What to change next.} Price the stock at 360 dollars (the grid becomes ten times finer relative to the price) and redraw the \omterm{def:qm:high-frequency-econometrics:signature}{signature plot}; let trades arrive in
Hawkes clusters (chapter~7) and see whether the kernel's bandwidth must change.
\end{tutorial}

\section{Build: the high-frequency estimators}\label{bld:qm:high-frequency-econometrics:hfvol}

\begin{build}
\bfield{Purpose} Every intraday volatility and correlation the miniature firm uses for quoting and risk is computed from ticks here, with the noise handled explicitly.
\bfield{Interface} \texttt{rv(logp, step)}, \texttt{rv\_subsampled}, \texttt{noise\_variance}, \texttt{tsrv(logp, k)}, \texttt{realised\_kernel(logp, H)},
\texttt{preaveraged(logp, kn)}, \texttt{bipower}, \texttt{optimal\_n(iq, noise\_var)}, \texttt{roll\_spread(price)}, \texttt{hayashi\_yoshida(t1, x1, t2, x2)},
\texttt{refresh\_times(*times)}, \texttt{previous\_tick(t, x, grid)}.
\bfield{Rules} No \omterm{def:qm:volatility-models:rv}{realised variance} is reported without its sampling interval; noise-robust \omterm{def:qm:estimation:estimator}{estimators} by default for anything sampled faster than a minute; covariances of
asynchronous assets by Hayashi--Yoshida or a synchronised noise-robust \omterm{def:qm:estimation:estimator}{estimator}, never by previous-tick returns on a fine grid.
\bfield{Acceptance tests} \texttt{code/firm/hfvol/tests/}: RV without noise equals IV; the noise bias $2n\omega^2$; two-scales, kernel and pre-averaging unbiased under iid noise;
\omterm{def:qm:high-frequency-econometrics:bipower}{bipower variation} ignores a jump; Hayashi--Yoshida unbiased under asynchronous Poisson trading and exact on a hand example; refresh times and previous-tick sampling;
\omterm{def:qm:high-frequency-econometrics:roll}{Roll's estimator} on pure bounce; the optimal $n$ minimises the approximate error.
\bfield{Stretch} The multivariate \omterm{def:qm:high-frequency-econometrics:kernel}{realised kernel}; jump tests based on the ratio of bipower to \omterm{def:qm:volatility-models:rv}{realised variance}; noise that depends on the trade sign and the spread.
\end{build}

\begin{omsources}
\item T.~G.~Andersen, T.~Bollerslev, F.~X.~Diebold and P.~Labys, ``Great realizations'', \emph{Risk}, March 2000, pp.~105--108.
\item R.~Roll, ``A simple implicit measure of the effective bid-ask spread in an efficient market'', \emph{Journal of Finance} 39, 1984.
\item T.~W.~Epps, ``Comovements in stock prices in the very short run'', \emph{Journal of the American Statistical Association} 74, 1979.
\item L.~Zhang, P.~A.~Mykland and Y.~Aït-Sahalia, ``A tale of two time scales'', \emph{Journal of the American Statistical Association} 100, 2005.
\item O.~E.~Barndorff-Nielsen, P.~R.~Hansen, A.~Lunde and N.~Shephard, ``Designing realized kernels to measure the ex post variation of equity prices in the presence of
noise'', \emph{Econometrica} 76, 2008.
\item J.~Jacod, Y.~Li, P.~A.~Mykland, M.~Podolskij and M.~Vetter, ``Microstructure noise in the continuous case: the pre-averaging approach'', \emph{Stochastic Processes and
their Applications} 119, 2009.
\item T.~Hayashi and N.~Yoshida, ``On covariance estimation of non-synchronously observed diffusion processes'', \emph{Bernoulli} 11, 2005.
\item F.~M.~Bandi and J.~R.~Russell, ``Microstructure noise, realized variance, and optimal sampling'', \emph{Review of Economic Studies} 75, 2008.
\item O.~E.~Barndorff-Nielsen and N.~Shephard, ``Power and bipower variation with stochastic volatility and jumps'', \emph{Journal of Financial Econometrics} 2, 2004.
\end{omsources}

\section{Exercises}

\begin{exercise}[$\star$]\label{exo:qm:high-frequency-econometrics:1}
A stock's daily \omterm{def:qm:volatility-models:rv}{integrated variance} is $(1\%)^2$ and its iid noise has a standard deviation of 2 basis points. What is the expected \omterm{def:qm:volatility-models:rv}{realised variance} from 23\,400 one-second
returns, and the implied daily volatility?
\end{exercise}

\begin{exercise}[$\star$]\label{exo:qm:high-frequency-econometrics:2}
Trade-price changes have first autocovariance $-0.0004$ (dollars squared). What spread does \omterm{def:qm:high-frequency-econometrics:roll}{Roll's estimator} imply?
\end{exercise}

\begin{exercise}[$\star$]\label{exo:qm:high-frequency-econometrics:3}
Why do mid quotes show less noise than trade prices in the \omterm{def:qm:high-frequency-econometrics:signature}{signature plot}, but still some?
\end{exercise}

\begin{exercise}[$\star\star$]\label{exo:qm:high-frequency-econometrics:4}
Compute the Bandi--Russell optimal number of returns and interval for a daily volatility of 1\% (constant) and noise of 2 basis points, over 23\,400 seconds.
\end{exercise}

\begin{exercise}[$\star\star$]\label{exo:qm:high-frequency-econometrics:5}
Show that with iid noise, observed high-frequency returns have first autocorrelation $-\omega^2/(\sigma^2\Delta + 2\omega^2)$, and evaluate it for the chapter's stock at one second.
\end{exercise}

\begin{exercise}[$\star\star$]\label{exo:qm:high-frequency-econometrics:6}
Two assets trade on average every 5 and 15 seconds. Roughly what share of one-second intervals contains a trade in both?
\end{exercise}

\begin{exercise}[$\star\star\star$]\label{exo:qm:high-frequency-econometrics:7}
\emph{Coding.} Draw the \omterm{def:qm:high-frequency-econometrics:signature}{signature plot} of the simulated stock from trades, mid quotes and efficient prices, and compute the root-mean-squared error of \omterm{def:qm:volatility-models:rv}{realised variance} at
each interval.
\end{exercise}

\begin{exercise}[$\star\star\star$]\label{exo:qm:high-frequency-econometrics:8}
\emph{Find the flaw.} ``To estimate the correlation of two stocks as precisely as possible, we used every one-second return of both over a year; the correlation is 0.12,
so the pair is useless for hedging.''
\end{exercise}

\section{Problem: The Volatility That Depended on the Clock}

\begin{problem}[{Weekend problem --- sampling a noisy price}]\label{pb:qm:high-frequency-econometrics:1}
A simulated stock has efficient volatility 25\% a year, trades once a second at the bid or ask of a one-cent grid around 36 dollars, over 23\,400-second days.

\medskip\noindent\textbf{Part I --- The \omterm{def:qm:high-frequency-econometrics:signature}{signature plot}.}
\begin{enumerate}
\item What volatility do one-second, five-second, one-minute and five-minute trade returns give?
\item And mid quotes at one second? The efficient price?
\item What is the noise standard deviation, estimated and true?
\item What does \omterm{def:qm:high-frequency-econometrics:roll}{Roll's estimator} say about the spread?
\item Why is the one-second variance 5.8 times too large?
\end{enumerate}

\medskip\noindent\textbf{Part II --- The optimal interval.}
\begin{enumerate}[resume]
\item What noise-to-signal ratio do the data imply?
\item What is the Bandi--Russell optimal number of returns and interval?
\item How does the simulated root-mean-squared error vary with the interval?
\item What does the optimum cost in data?
\item How would the optimum change on a day twice as volatile?
\end{enumerate}

\medskip\noindent\textbf{Part III --- Beyond sampling.}
\begin{enumerate}[resume]
\item What do two-scales RV, pre-averaging and the \omterm{def:qm:high-frequency-econometrics:kernel}{realised kernel} give, and how precisely?
\item What is the \omterm{def:qm:high-frequency-econometrics:epps}{Epps effect} for two stocks trading every 5 and 15 seconds?
\item What do Hayashi--Yoshida and \omterm{def:qm:high-frequency-econometrics:hy}{refresh-time sampling} give?
\item What does a 1\% jump do to realised and \omterm{def:qm:high-frequency-econometrics:bipower}{bipower variation} at five minutes?
\item Why can't jumps and noise be handled separately?
\end{enumerate}

\medskip\noindent\textbf{Part IV --- Judgement.}
\begin{enumerate}[resume]
\item Which intraday volatility should a market maker's quoting engine use?
\item Which correlation should its hedging use?
\item What would you check on real data that the simulation assumes away?
\item State the \emph{named result}: the optimal sampling interval and the bias of one-second \omterm{def:qm:volatility-models:rv}{realised variance} it avoids.
\item In one sentence: why does more data make \omterm{def:qm:volatility-models:rv}{realised variance} worse?
\end{enumerate}
\end{problem}

\section{Interview questions}

\begin{interviewq}[$\star$ \iqroles{researcher, trader}]\label{iq:qm:high-frequency-econometrics:1}
Why does volatility measured from tick data come out higher than from five-minute data?
\end{interviewq}

\begin{interviewq}[$\star\star$ \iqroles{researcher, developer}]\label{iq:qm:high-frequency-econometrics:2}
How would you estimate the spread of a stock from trade prices alone?
\end{interviewq}

\begin{interviewq}[$\star\star$ \iqroles{researcher}]\label{iq:qm:high-frequency-econometrics:3}
What is the optimal sampling frequency for \omterm{def:qm:volatility-models:rv}{realised variance}, and what does it depend on?
\end{interviewq}

\begin{interviewq}[$\star\star$ \iqroles{researcher, trader}]\label{iq:qm:high-frequency-econometrics:4}
Your realised correlation between two liquid stocks falls from 0.6 to 0.1 when you move from five-minute to one-second data. Why, and how do you fix it?
\end{interviewq}

\begin{interviewq}[$\star\star$ \iqroles{researcher}]\label{iq:qm:high-frequency-econometrics:5}
How do you separate jumps from diffusive volatility in intraday data?
\end{interviewq}

\begin{interviewq}[$\star\star\star$ \iqroles{researcher}]\label{iq:qm:high-frequency-econometrics:6}
Explain why a \omterm{def:qm:high-frequency-econometrics:kernel}{realised kernel} removes \omterm{def:qm:high-frequency-econometrics:noise}{microstructure noise}, in terms of autocovariances.
\end{interviewq}
